\section{Discussion}
\label{sec:discussion}

\paragraph{The signal is still there, but it is thin.}
A gross Sharpe ratio between 1.2 and 1.5 for a market-neutral strategy is a strong result in
statistical terms. Residual mean reversion in large US stocks has not disappeared since AL's
sample. Economically, however, the edge is about 2.5--5~bps per unit traded, comparable to the
cost of trading a large-cap stock. Whether the strategy is viable therefore depends almost
entirely on execution costs, and a cost assumption of 5~bps versus 2~bps would turn a losing
strategy into a profitable one. Several explanations are plausible for why the edge is smaller
than in AL's 1997--2007 sample: more capital competing for the same signal \citep{khandani2007},
narrower spreads and faster price discovery after decimalisation and the rise of electronic
market making, and a universe restricted to the 500 largest and most liquid US stocks.

\paragraph{Hedging is the hidden cost.}
In the PCA formulation, factor hedges are dense portfolios across every stock. Hedging is
therefore not a small adjustment: it doubles the trading volume of the strategy. The flat 5~bps
cost model treats a hedge trade in the same way as a signal trade. In practice, the systematic
exposure could be hedged far more cheaply with index futures or sector ETFs, whose trading costs
are a fraction of a basis point to about one basis point. This suggests an important extension:
hedge with liquid instruments (AL's ETF variant) and charge each leg its own realistic cost.

\paragraph{Timing matters as much as costs.}
The execution-lag results (\cref{sec:results:lag}) show that conclusions drawn from AL's
convention of trading at the signal's close are optimistic. Short-horizon reversal is exactly the
kind of signal that bid--ask bounce inflates: a closing price that prints at the bid looks like a
loss that reverses the next day, even if nothing happened. A realistic evaluation therefore needs
both a delay and a cost model, and the two interact. The strategies most robust to the delay
(fewer factors, the no-trade band, the deep model) are also the ones that trade least.

\paragraph{What the deep model adds, and what it does not.}
Training on the net Sharpe ratio, with costs charged on the implied stock trades, produced a model
that trades far less and keeps most of its edge under a one-day lag. Its break-even cost of about
4~bps is the highest of any strategy evaluated at lag 1 without in-sample selection. Two caveats
qualify this. First, the model learned a nearly long-only residual book, so about half of its
lagged edge is an equal-weighted long position in residuals: a systematic tilt that a simple,
cheap portfolio could capture, not statistical arbitrage. Constraining the residual book to be
balanced (for example by demeaning scores before \cref{eq:deep-alloc}) would isolate the reversal
component; this is left for future work, since it would add trials. Second, its validation
performance is unstable across years, so the five-year average hides two losing years.

\paragraph{Fewer factors, chosen by random-matrix theory.}
The peak of the gross Sharpe ratio at $m=8$, matching the number of eigenvalues above the
Marchenko--Pastur edge, suggests choosing $m$ adaptively each day as the number of
eigenvalues that are statistically distinguishable from noise. This is consistent with the
finding that the explained variance varies strongly over time (\cref{fig:explained}). It is
only suggestive so far, since the edge count was examined on a single date.

\paragraph{Extreme s-scores are not better signals.}
The finding that wider entry thresholds lower the profit per trade suggests that very large
residual moves are more often driven by information, such as earnings or news, than by
temporary liquidity imbalances. Such moves are less likely to revert. Volume adjustment was
meant to address exactly this, but it did not help in our sample. Conditioning on scheduled
events such as earnings announcements is a more direct alternative.

\subsection{Limitations}
\label{sec:limitations}

\begin{itemize}
  \item \textbf{Survivorship bias.} The universe is the 2026 S\&P~500 membership applied to
        2010--2025 (\cref{sec:data}).
  \item \textbf{Cost model.} A flat proportional cost, with no market impact, no short-borrow
        fees and no distinction between liquid and illiquid names or between stock and hedge
        trades.
  \item \textbf{Execution.} Only closing prices are available. A one-day lag is tested, but
        execution at the next open or within the day is not.
  \item \textbf{Accounting.} Fixed capital, no compounding, no margin or financing costs, and
        missing next-day returns treated as zero.
  \item \textbf{Multiple testing.} \NVariants{} rule-based configurations and the deep model at two
        execution lags have been evaluated in-sample. The best in-sample Sharpe ratios are biased
        upwards and must be deflated.
  \item \textbf{Single market and period.} Results are for large US stocks in 2011--2018 only.
\end{itemize}
